\documentclass[a4paper,12pt]{article}

% --------------------------------------------------
% PAQUETES
% --------------------------------------------------
\usepackage[utf8]{inputenc}
\usepackage[T1]{fontenc}
\usepackage[spanish,es-nodecimaldot]{babel}
\usepackage{lmodern}

\usepackage[a4paper,margin=2.6cm]{geometry}
\usepackage{amsmath,amssymb,amsthm,mathtools,physics}
\usepackage{array,booktabs,multirow}
\usepackage{float}
\usepackage{caption}
\usepackage{tikz}
\usepackage{pgfplots}
\usepackage{hyperref}
\usepackage{fancyhdr}
\usepackage{mdframed}
\usepackage{titlesec}
\usepackage{enumitem}
\usepackage{graphicx}
\usepackage{xcolor}
\usepackage{setspace}

\pgfplotsset{compat=1.18}
\usetikzlibrary{arrows.meta, decorations.pathreplacing, calc}

% --------------------------------------------------
% HIPERVÍNCULOS
% --------------------------------------------------
\hypersetup{
    colorlinks=true,
    linkcolor=blue!50!black,
    urlcolor=blue!50!black,
    citecolor=blue!50!black,
    pdftitle={Métodos Numéricos - Apuntes Unificados},
    pdfauthor={Luis López},
    pdfsubject={Métodos Numéricos},
    pdfkeywords={Métodos numéricos, ecuaciones no lineales,
                 interpolación, mínimos cuadrados}
}

% --------------------------------------------------
% ESPACIADO
% --------------------------------------------------
\setlength{\parskip}{0.6em}
\setlength{\parindent}{0pt}
\onehalfspacing

% --------------------------------------------------
% MACROS
% --------------------------------------------------
\newcommand{\R}{\mathbb{R}}
\newcommand{\C}{\mathbb{C}}
\newcommand{\N}{\mathbb{N}}
\newcommand{\Z}{\mathbb{Z}}
\newcommand{\eps}{\varepsilon}
\renewcommand{\vec}[1]{\mathbf{#1}}
\newcommand{\ve}[1]{\mathbf{#1}}
\newcommand{\uvec}[1]{\hat{\mathbf{#1}}}

\newcommand{\EN}{E_N}
\newcommand{\EA}{E_a}
\newcommand{\ER}{E_r}
\newcommand{\Tol}{T_s}
\newcommand{\fl}{\mathrm{fl}}

% --------------------------------------------------
% ESTILO DE SECCIONES
% --------------------------------------------------
\titleformat{\section}
{\normalfont\Large\bfseries}
{\thesection.}{0.6em}{}
[\vspace{0.25em}\titlerule]

\titleformat{\subsection}
{\normalfont\large\bfseries}
{\thesubsection.}{0.5em}{}

\titleformat{\subsubsection}
{\normalfont\normalsize\bfseries}
{\thesubsubsection.}{0.5em}{}

% --------------------------------------------------
% CABECERA Y PIE
% --------------------------------------------------
\pagestyle{fancy}
\fancyhf{}
\fancyhead[L]{\nouppercase{\leftmark}}
\fancyhead[R]{\thepage}
\renewcommand{\headrulewidth}{0.4pt}

% --------------------------------------------------
% ENTORNOS
% --------------------------------------------------
\theoremstyle{definition}
\newtheorem{definicion}{Definición}[section]
\newtheorem{ejemplo}{Ejemplo}[section]
\newtheorem{observacion}{Observación}[section]
\newtheorem{nota}{Nota}[section]

\theoremstyle{plain}
\newtheorem{teorema}{Teorema}[section]
\newtheorem{proposicion}{Proposición}[section]
\newtheorem{lema}{Lema}[section]
\newtheorem{corolario}{Corolario}[section]

% --------------------------------------------------
% CAJAS
% --------------------------------------------------
\mdfdefinestyle{resultado}{
    linecolor=black,
    linewidth=0.8pt,
    roundcorner=4pt,
    backgroundcolor=gray!8,
    innertopmargin=0.8em,
    innerbottommargin=0.8em,
    innerleftmargin=0.8em,
    innerrightmargin=0.8em
}

\newmdenv[style=resultado]{cajaresultado}

% --------------------------------------------------
% DOCUMENTO
% --------------------------------------------------
\begin{document}

% --------------------------------------------------
% PORTADA
% --------------------------------------------------
\begin{titlepage}
    \centering
    \vspace*{1.5cm}

    {\Huge \textbf{Métodos Numéricos}\par}
    \vspace{0.3cm}
    {\Large Curso 2025--2026\par}
    \vspace{0.3cm}
    {\Large Apuntes Unificados\par}
    \vspace{0.3cm}
    {\large Teoría, demostraciones, desarrollos y formulario\par}
    \vspace{0.7cm}
    {\large Temas 1--3: error numérico, ecuaciones no lineales,
             interpolación y aproximación\par}

    \vspace{1.2cm}

    \begin{tikzpicture}[scale=1.05]
        % Ejes
        \draw[->,thick] (-5.2,0) -- (5.3,0) node[right] {$x$};
        \draw[->,thick] (0,-3.2) -- (0,4.1) node[above] {$y$};

        % Curva exacta
        \draw[blue!70!black, very thick, domain=-4.2:4.1, samples=200]
            plot (\x,{0.12*(\x)^3 - 0.75*(\x) + 0.6});

        % Aproximaciones de Taylor (visual)
        \draw[red!70!black, thick] (-4.3,-2.2) -- (4.4,2.7);

        % Puntos de interpolación
        \filldraw[black] (-2.5,1.1) circle (2pt);
        \filldraw[black] (-1.4,0.6) circle (2pt);
        \filldraw[black] (0.8,-0.1) circle (2pt);
        \filldraw[black] (2.0,-0.35) circle (2pt);
        \filldraw[black] (3.0,1.4) circle (2pt);

        % Etiquetas
        \node[blue!70!black] at (3.7,3.0) {$f(x)$};
        \node[red!70!black] at (3.9,2.1) {aprox.};
        \node at (-3.7,-2.6)
              {\small error, raíces, ajuste, interpolación};
    \end{tikzpicture}

    \vfill

    {\large Autor: Luis López\par}
    {\large Titulación: Grado en Física\par}
    {\large Fecha: Marzo 2026\par}

    \vspace{0.6cm}
\end{titlepage}

\tableofcontents
\newpage

% ==================================================
% TEMA 1
% ==================================================
\section{Tema 1: Error por truncamiento y por redondeo}

\subsection{Introducción}

En cálculo numérico no basta con producir una aproximación; también es esencial
entender \emph{qué error se comete}, \emph{de dónde procede} y
\emph{cómo se controla}. El error numérico aparece de forma natural porque los
algoritmos trabajan con representaciones finitas y porque, en muchas ocasiones,
se sustituyen expresiones exactas por aproximaciones más manejables.

En este tema se estudian dos mecanismos fundamentales de error:

\begin{itemize}[leftmargin=1.5em]
    \item el \textbf{error por truncamiento}, que surge al sustituir una expresión
          exacta por otra incompleta, por ejemplo al cortar una serie de Taylor o
          Maclaurin;
    \item el \textbf{error por redondeo}, que aparece al representar números reales
          con un número finito de cifras.
\end{itemize}

Estos dos tipos de error son la base conceptual del análisis de la estabilidad y
precisión de cualquier método numérico. El análisis riguroso de ambos es
imprescindible: un algoritmo puede ser matemáticamente correcto y, sin embargo,
dar resultados completamente erróneos si no se controlan los errores numéricos.

\subsection{Conceptos básicos de error}

\begin{definicion}[Error absoluto]
Sea $x$ el valor exacto y $\tilde{x}$ una aproximación. El \textbf{error absoluto}
se define por
\[
    E_a = |x - \tilde{x}|.
\]
\end{definicion}

\begin{definicion}[Error relativo]
El \textbf{error relativo} asociado a $\tilde{x}$ es
\[
    E_r = \frac{|x - \tilde{x}|}{|x|},
    \qquad x \neq 0.
\]
Frecuentemente se expresa en porcentaje:
\[
    E_r(\%) = \frac{|x - \tilde{x}|}{|x|} \cdot 100.
\]
\end{definicion}

\begin{observacion}
El error absoluto mide la desviación en las mismas unidades de la magnitud,
mientras que el error relativo mide la importancia del error con respecto al
tamaño del valor exacto. Para comparar aproximaciones de magnitudes distintas,
el error relativo es la medida adecuada.
\end{observacion}

\begin{definicion}[Cifras correctas y decimales correctos]
Se dice que una aproximación $\tilde{x}$ posee $k$ \textbf{decimales correctos}
si el error relativo satisface, de forma aproximada,
\[
    E_r < \frac{1}{2}\,10^{-k}.
\]
Una fórmula muy utilizada para estimar el número de decimales correctos es
\[
    k \approx \left\lfloor -\log_{10}(2E_r)\right\rfloor.
\]
\end{definicion}

\begin{cajaresultado}
\textbf{Fórmulas fundamentales de error:}
\[
    E_a = |x - \tilde{x}|,
    \qquad
    E_r = \frac{|x - \tilde{x}|}{|x|},
    \qquad
    k \approx \left\lfloor -\log_{10}(2E_r) \right\rfloor.
\]
\end{cajaresultado}

\subsection{Error por truncamiento}

\subsubsection{Idea general}

\begin{definicion}[Error por truncamiento]
El \textbf{error por truncamiento} es el error cometido al reemplazar un proceso
infinito por uno finito. Por ejemplo, si una función se desarrolla en serie de
Taylor y se conservan sólo los primeros $n+1$ términos, la diferencia entre la
función exacta y el polinomio truncado es el error de truncamiento.
\end{definicion}

La herramienta central para este estudio es la serie de Taylor.

\subsubsection{Serie de Taylor y resto de Lagrange}

\begin{teorema}[Serie de Taylor con resto de Lagrange]
Sea $f$ una función con derivadas hasta orden $n+1$ en un entorno de $a$.
Entonces
\[
    f(x) = \sum_{k=0}^{n} \frac{f^{(k)}(a)}{k!}(x-a)^k + R_n(x),
\]
donde el \textbf{resto de Lagrange} es
\[
    R_n(x) = \frac{f^{(n+1)}(\xi)}{(n+1)!}(x-a)^{n+1},
\]
para algún $\xi$ situado estrictamente entre $a$ y $x$.
\end{teorema}

\begin{demostracion}
Se aplica el teorema del valor medio generalizado de forma recursiva.
Definimos el resto exacto $R_n(x) = f(x) - T_n(x)$, donde
$T_n(x) = \sum_{k=0}^{n} \frac{f^{(k)}(a)}{k!}(x-a)^k$.
Consideremos la función auxiliar
\[
    G(t) = f(x) - \sum_{k=0}^{n}\frac{f^{(k)}(t)}{k!}(x-t)^k
           - \frac{R_n(x)}{(x-a)^{n+1}}(x-t)^{n+1}.
\]
Se verifica $G(x) = 0$ y $G(a) = 0$. Por el teorema de Rolle existe
$\xi \in (a,x)$ con $G'(\xi) = 0$. Calculando $G'(t)$ y simplificando
(los términos del sumatorio se cancelan telescópicamente), se obtiene
\[
    G'(\xi) = -\frac{f^{(n+1)}(\xi)}{n!}(x-\xi)^n
              + \frac{(n+1)R_n(x)}{(x-a)^{n+1}}(x-\xi)^n = 0,
\]
de donde
\[
    R_n(x) = \frac{f^{(n+1)}(\xi)}{(n+1)!}(x-a)^{n+1}. \qedhere
\]
\end{demostracion}

\begin{cajaresultado}
\textbf{Error de truncamiento de orden $n$:}
\[
    f(x) = T_n(x) + R_n(x),
    \qquad
    E_{\mathrm{trunc}} = |R_n(x)| = \left|\frac{f^{(n+1)}(\xi)}{(n+1)!}
    (x-a)^{n+1}\right|.
\]
El error de truncamiento \emph{disminuye} al aumentar $n$ (siempre que la serie
converja) y \emph{aumenta} cuando el punto de evaluación $x$ se aleja del centro
$a$.
\end{cajaresultado}

\begin{observacion}
El punto $\xi$ que aparece en el resto de Lagrange es \emph{desconocido}; sólo
sabemos que pertenece al intervalo $(a,x)$. En la práctica, el resto se usa para
acotar el error, tomando el máximo de $|f^{(n+1)}|$ en ese intervalo.
\end{observacion}

\subsubsection{Serie de Maclaurin}

Cuando el desarrollo se centra en $a = 0$, se obtiene la \textbf{serie de
Maclaurin}:
\[
    f(x) = f(0) + f'(0)x + \frac{f''(0)}{2!}x^2 + \cdots
           + \frac{f^{(n)}(0)}{n!}x^n + R_n(x).
\]

\begin{proposicion}[Serie de Maclaurin de $e^x$]
La serie de Maclaurin de la función exponencial es
\[
    e^x = \sum_{k=0}^{\infty} \frac{x^k}{k!}
        = 1 + x + \frac{x^2}{2!} + \frac{x^3}{3!} + \cdots
\]
Esta serie converge para todo $x \in \R$.
\end{proposicion}

\begin{demostracion}
Toda derivada de $e^x$ coincide con la propia función. En particular,
$f^{(k)}(0) = e^0 = 1$ para todo $k \in \N$. Sustituyendo en la fórmula
general de Maclaurin:
\[
    e^x = \sum_{k=0}^{n} \frac{x^k}{k!} + R_n(x),
    \qquad
    R_n(x) = \frac{e^\xi}{(n+1)!} x^{n+1},
    \quad \xi \in (0, x).
\]
Para cualquier $x$ fijo, $|R_n(x)| \leq \frac{e^{|x|}}{(n+1)!}|x|^{n+1}
\to 0$ cuando $n \to \infty$, ya que el factorial crece más rápido que
cualquier potencia. Por tanto la serie converge para todo $x \in \R$.
\end{demostracion}

\subsubsection{Aproximaciones sucesivas de Taylor}

Sea
\[
    T_0(x) = f(0), \quad
    T_1(x) = f(0) + f'(0)x, \quad
    T_2(x) = f(0) + f'(0)x + \frac{f''(0)}{2!}x^2,
\]
y en general
\[
    T_n(x) = \sum_{k=0}^{n} \frac{f^{(k)}(0)}{k!} x^k.
\]

Cuanto mayor es el orden $n$, menor suele ser el error por truncamiento,
siempre que la serie converja adecuadamente en el punto de evaluación.
La figura siguiente ilustra cómo los polinomios de Taylor de orden creciente
aproximan mejor a $e^x$ en un entorno del origen.

\begin{figure}[H]
\centering
\begin{tikzpicture}
\begin{axis}[
    width=11cm, height=7cm,
    xmin=-2.5, xmax=2.5,
    ymin=-1, ymax=8,
    xlabel={$x$}, ylabel={$y$},
    grid=major,
    grid style={gray!25},
    legend pos=north west,
    legend style={font=\small},
    tick label style={font=\small}
]
% Función exacta
\addplot[blue!70!black, very thick, domain=-2.5:2.5, samples=200]
    {exp(x)};
\addlegendentry{$e^x$ (exacta)}
% T0
\addplot[red, thick, dashed, domain=-2.5:2.5]
    {1};
\addlegendentry{$T_0 = 1$}
% T1
\addplot[orange!80!black, thick, domain=-2.5:2.5]
    {1 + x};
\addlegendentry{$T_1 = 1 + x$}
% T2
\addplot[green!60!black, thick, domain=-2.5:2.5]
    {1 + x + x^2/2};
\addlegendentry{$T_2$}
% T3
\addplot[purple!80!black, thick, domain=-2.5:2.5]
    {1 + x + x^2/2 + x^3/6};
\addlegendentry{$T_3$}
% T4
\addplot[cyan!70!black, thick, domain=-2.5:2.5]
    {1 + x + x^2/2 + x^3/6 + x^4/24};
\addlegendentry{$T_4$}
\end{axis}
\end{tikzpicture}
\caption{Aproximaciones sucesivas de Taylor de $e^x$ en torno a $a=0$.
A mayor orden, mejor aproximación en el intervalo mostrado.}
\end{figure}

\subsubsection{Ejemplo guiado: error de truncamiento con datos tabulados}

\begin{ejemplo}
Supóngase que se conocen los datos
\[
    f(0) = 1.2, \quad f'(0) = -0.25, \quad f''(0) = -1.0, \quad
    f^{(3)}(0) = -0.9, \quad f^{(4)}(0) = -2.4,
\]
y se desea aproximar $f(1)$, sabiendo que el valor exacto es $f(1) = 0.2$.

El desarrollo de Taylor en torno a $a = 0$ es
\[
    f(x) \approx 1.2 - 0.25x - \frac{1.0}{2}x^2
                - \frac{0.9}{6}x^3 - \frac{2.4}{24}x^4
        = 1.2 - 0.25x - 0.5x^2 - 0.15x^3 - 0.1x^4.
\]

Evaluando en $x = 1$ con órdenes crecientes:

\medskip
\begin{center}
\begin{tabular}{clcc}
\toprule
Orden & Polinomio $T_n(1)$ & $E_a = |0.2 - T_n(1)|$ & $E_r$ \\
\midrule
0 & $T_0(1) = 1.2$ & $1.0000$ & $500.00\%$ \\
1 & $T_1(1) = 1.2 - 0.25 = 0.95$ & $0.7500$ & $375.00\%$ \\
2 & $T_2(1) = 0.95 - 0.50 = 0.45$ & $0.2500$ & $125.00\%$ \\
3 & $T_3(1) = 0.45 - 0.15 = 0.30$ & $0.1000$ & $50.00\%$ \\
4 & $T_4(1) = 0.30 - 0.10 = 0.20$ & $0$ & $0\%$ \\
\bottomrule
\end{tabular}
\end{center}

En este caso el polinomio de orden 4 reproduce \emph{exactamente} el valor en
$x = 1$. El término residual de Lagrange es
\[
    R_4(1) = \frac{f^{(5)}(\xi)}{5!} \cdot 1^5.
\]
Si la función es un polinomio de grado 4, entonces $f^{(5)} \equiv 0$ y
$R_4 = 0$, lo que explica el error nulo.
\end{ejemplo}

\begin{observacion}
El hecho de que en un ejemplo concreto el error de truncamiento sea nulo no
implica que esto ocurra en general. Sólo indica que, para esos datos, el
polinomio truncado de orden 4 coincide con la función en el punto considerado.
Si $f$ tuviese derivadas de orden superior no nulas, la aproximación de orden 4
en otros puntos tendría error no nulo.
\end{observacion}

\subsubsection{Ejemplo detallado: aproximación de $e^{0.5}$ con tolerancia}

\begin{ejemplo}
\label{ej:exp05}
Aproximar $e^{0.5}$ mediante su serie de Maclaurin hasta lograr una tolerancia
del orden de $10^{-4}$. Se toma como valor exacto $e^{0.5} = 1.648721271$.

Partimos de la serie:
\[
    e^x = \sum_{k=0}^{\infty} \frac{x^k}{k!},
    \qquad
    e^{0.5} = \sum_{k=0}^{\infty} \frac{(0.5)^k}{k!}.
\]

Para cada orden $n$ calculamos la aproximación $T_n$, el error absoluto
$E_a = |1.648721271 - T_n|$ y el error relativo $E_r = E_a / 1.648721271$.
También estimamos el número de decimales correctos mediante
$k = \lfloor -\log_{10}(2E_r)\rfloor$.

\medskip
\textbf{Cálculo paso a paso:}

\begin{align*}
T_0 &= 1, \\
T_1 &= 1 + 0.5 = 1.5, \\
T_2 &= 1.5 + \frac{(0.5)^2}{2!} = 1.5 + 0.125 = 1.625, \\
T_3 &= 1.625 + \frac{(0.5)^3}{3!} = 1.625 + 0.020\overline{8}
     = 1.645833\overline{3}, \\
T_4 &= 1.645833\overline{3} + \frac{(0.5)^4}{4!}
     = 1.645833\overline{3} + 0.002604167 = 1.6484375, \\
T_5 &= 1.6484375 + \frac{(0.5)^5}{5!}
     = 1.6484375 + 0.000260417 = 1.648697917, \\
T_6 &= 1.648697917 + \frac{(0.5)^6}{6!}
     = 1.648697917 + 0.0000217\overline{6} = 1.648719618.
\end{align*}

\medskip
\begin{center}
\begin{tabular}{ccccc}
\toprule
Orden $n$ & $T_n$ & $E_a$ & $E_r$ & Decimales correctos $k$ \\
\midrule
0 & 1.000000 & $6.49 \times 10^{-1}$ & $3.94 \times 10^{-1}$ & 0 \\
1 & 1.500000 & $1.49 \times 10^{-1}$ & $9.02 \times 10^{-2}$ & 0 \\
2 & 1.625000 & $2.37 \times 10^{-2}$ & $1.44 \times 10^{-2}$ & 1 \\
3 & 1.645833 & $2.89 \times 10^{-3}$ & $1.75 \times 10^{-3}$ & 2 \\
4 & 1.648438 & $2.84 \times 10^{-4}$ & $1.72 \times 10^{-4}$ & 3 \\
\textbf{5} & \textbf{1.648698} & $\mathbf{2.34 \times 10^{-5}}$
           & $\mathbf{1.42 \times 10^{-5}}$ & \textbf{4} \\
6 & 1.648720 & $1.65 \times 10^{-6}$ & $1.00 \times 10^{-6}$ & 5 \\
\bottomrule
\end{tabular}
\end{center}

\medskip
\textbf{Verificación de la tolerancia con la fórmula de cifras correctas:}

Para $n = 5$: $E_r = 1.42 \times 10^{-5}$, luego
\[
    k = \left\lfloor -\log_{10}(2 \times 1.42 \times 10^{-5}) \right\rfloor
      = \left\lfloor -\log_{10}(2.84 \times 10^{-5}) \right\rfloor
      = \left\lfloor 4.547 \right\rfloor = 4.
\]
Esto confirma que con $n = 5$ se obtienen 4 decimales correctos, suficiente
para satisfacer la tolerancia $10^{-4}$.
\end{ejemplo}

\begin{cajaresultado}
Para aproximar $e^{0.5}$ con tolerancia del orden de $10^{-4}$ basta tomar
el polinomio de Maclaurin de \textbf{orden 5}:
\[
    e^{0.5} \approx 1.6487.
\]
A partir de $n = 5$, $E_a < 10^{-4}$ y se tienen al menos 4 decimales
correctos.
\end{cajaresultado}

\begin{figure}[H]
\centering
\begin{tikzpicture}
\begin{axis}[
    width=11cm, height=6.5cm,
    xlabel={Orden $n$},
    ylabel={$\log_{10}(E_a)$},
    xtick={0,1,2,3,4,5,6},
    grid=major,
    grid style={gray!25},
    ymin=-7, ymax=0,
    legend pos=north east,
    tick label style={font=\small}
]
\addplot[blue!70!black, very thick, mark=*, mark size=2.5pt]
    coordinates {
        (0,-0.188)
        (1,-0.827)
        (2,-1.625)
        (3,-2.539)
        (4,-3.547)
        (5,-4.631)
        (6,-5.783)
    };
\addlegendentry{$\log_{10}(E_a)$}
% Línea de tolerancia 1e-4
\addplot[red!70!black, dashed, thick, domain=0:6] {-4};
\addlegendentry{Tolerancia $10^{-4}$}
\end{axis}
\end{tikzpicture}
\caption{Convergencia del error absoluto al aproximar $e^{0.5}$ con la serie
de Maclaurin. A partir del orden $n=5$ se cruza la tolerancia $10^{-4}$.}
\end{figure}

\subsubsection{Ejercicio resuelto: serie de Maclaurin de $\dfrac{1}{1+x^2}$}

\begin{ejemplo}
Desarrollar $f(x) = \dfrac{1}{1+x^2}$ en serie de Maclaurin y determinar el
radio de convergencia.

\medskip
\textbf{Método 1 — Series geométricas (más elegante):}

Recordamos la serie geométrica:
\[
    \frac{1}{1-u} = \sum_{k=0}^{\infty} u^k = 1 + u + u^2 + u^3 + \cdots,
    \qquad |u| < 1.
\]
Identificamos $u = -x^2$:
\[
    \frac{1}{1+x^2} = \frac{1}{1-(-x^2)}
    = \sum_{k=0}^{\infty} (-x^2)^k
    = \sum_{k=0}^{\infty} (-1)^k x^{2k},
    \qquad |{-x^2}| < 1 \Rightarrow |x| < 1.
\]

Explícitamente:
\[
    \frac{1}{1+x^2} = 1 - x^2 + x^4 - x^6 + x^8 - \cdots
\]

\medskip
\textbf{Radio de convergencia:} La serie converge para $|x| < 1$.

\medskip
\textbf{Método 2 — Derivadas sucesivas (verificación directa):}

Calculamos las derivadas en $x = 0$:
\begin{align*}
f(x) &= (1+x^2)^{-1}
        &\Rightarrow\quad f(0) &= 1, \\
f'(x) &= -\frac{2x}{(1+x^2)^2}
         &\Rightarrow\quad f'(0) &= 0, \\
f''(x) &= \frac{-2(1+x^2)^2 + 8x^2(1+x^2)}{(1+x^2)^4}
          = \frac{6x^2 - 2}{(1+x^2)^3}
          &\Rightarrow\quad f''(0) &= -2, \\
f'''(x) &= \frac{24x - 24x^3}{(1+x^2)^4}
           &\Rightarrow\quad f'''(0) &= 0, \\
f^{(4)}(x)\big|_{x=0} &= 24.
\end{align*}

Aplicando la fórmula de Maclaurin $\sum_{k=0}^{n}
\frac{f^{(k)}(0)}{k!}x^k$:
\[
    \frac{1}{1+x^2} = 1 + 0 \cdot x + \frac{-2}{2!}x^2
                    + 0 \cdot x^3 + \frac{24}{4!}x^4 + \cdots
    = 1 - x^2 + x^4 - x^6 + \cdots \checkmark
\]

Ambos métodos son consistentes. Los coeficientes impares son nulos porque
$f(x) = f(-x)$ (función par).

\begin{cajaresultado}
\[
    \frac{1}{1+x^2} = \sum_{k=0}^{\infty} (-1)^k x^{2k}
    = 1 - x^2 + x^4 - x^6 + x^8 - \cdots,
    \qquad |x| < 1.
\]
Integrando término a término se obtiene la serie de $\arctan x$:
\[
    \arctan x = \sum_{k=0}^{\infty} \frac{(-1)^k x^{2k+1}}{2k+1}
    = x - \frac{x^3}{3} + \frac{x^5}{5} - \cdots,
    \qquad |x| \leq 1.
\]
\end{cajaresultado}
\end{ejemplo}

\subsection{Error por redondeo}

\subsubsection{Definición y origen}

\begin{definicion}[Error por redondeo]
El \textbf{error por redondeo} es el que aparece cuando un número real se
representa mediante un número finito de cifras. Si $x$ es el número exacto y
$\fl(x)$ su representación en máquina (punto flotante), entonces
\[
    E_{\mathrm{red}} = |x - \fl(x)|.
\]
\end{definicion}

En la práctica, el redondeo es inevitable porque los sistemas de cómputo sólo
pueden almacenar un número finito de dígitos. La norma IEEE 754 de doble
precisión usa 64 bits: 1 de signo, 11 de exponente y 52 de mantisa, lo que
proporciona aproximadamente 15--16 dígitos decimales significativos y un épsilon
de máquina $\epsilon_{\mathrm{máq}} \approx 2.22 \times 10^{-16}$.

\begin{definicion}[Épsilon de máquina]
El \textbf{épsilon de máquina} $\epsilon_{\mathrm{máq}}$ es el número positivo
más pequeño tal que
\[
    \fl(1 + \epsilon_{\mathrm{máq}}) > 1.
\]
Representa la precisión relativa de la aritmética de punto flotante: cualquier
número real $x$ se representa con error relativo acotado por
$\epsilon_{\mathrm{máq}}/2$.
\end{definicion}

\subsubsection{Truncar versus redondear}

Hay dos operaciones distintas para reducir el número de cifras de un número:

\begin{itemize}[leftmargin=1.5em]
    \item \textbf{Truncar (corte):} se eliminan los dígitos sobrantes sin
          modificar el último conservado. Por ejemplo, $\pi = 3.14159\ldots$
          truncado a 4 decimales es $3.1415$.
    \item \textbf{Redondear:} se modifica el último dígito conservado según los
          descartados, aplicando la regla de redondeo.
\end{itemize}

\begin{definicion}[Regla de redondeo simétrico]
Sea un número decimal y supóngase que se desea conservar cierto número de
cifras:
\begin{itemize}[leftmargin=1.5em]
    \item Si el primer dígito descartado es \emph{mayor} que 5, se incrementa
          en una unidad el último dígito conservado.
    \item Si el primer dígito descartado es \emph{menor} que 5, el último
          dígito conservado permanece igual.
    \item Si el primer dígito descartado es exactamente 5 (seguido de ceros),
          se aplica el \textbf{redondeo al par} (\emph{banker's rounding}): el
          último dígito conservado se incrementa sólo si es impar.
\end{itemize}
\end{definicion}

\begin{observacion}
El redondeo al par es la convención estándar en la norma IEEE 754 y reduce el
sesgo estadístico acumulado cuando se realizan muchas operaciones con números
que terminan exactamente en 5.
\end{observacion}

\subsubsection{Ejemplos de redondeo y truncamiento}

\begin{ejemplo}
Redondear a dos decimales:
\[
    5.6723 \mapsto 5.67
    \quad \text{(tercer decimal = 2 < 5, se deja igual)},
\]
\[
    10.406 \mapsto 10.41
    \quad \text{(tercer decimal = 6 > 5, se incrementa)}.
\]
\end{ejemplo}

\begin{ejemplo}
Redondear a un decimal (casos con dígito = 5):
\[
    7.3500 \mapsto 7.4
    \quad \text{(último conservado = 3, impar $\Rightarrow$ se incrementa)},
\]
\[
    1.2500 \mapsto 1.2
    \quad \text{(último conservado = 2, par $\Rightarrow$ se deja igual)}.
\]
\end{ejemplo}

\begin{ejemplo}
Redondear a tres decimales:
\[
    88.21650 \mapsto 88.216
    \quad \text{(cuarto decimal = 5, último conservado = 6, par
                $\Rightarrow$ se deja igual)}.
\]
Truncar a tres decimales:
\[
    88.21650 \mapsto 88.216 \quad \text{(coincide en este caso)},
    \qquad
    \pi = 3.14159\ldots \mapsto 3.141 \quad \text{(truncado)}.
\]
\end{ejemplo}

\begin{ejemplo}
Representación de $\pi$ en máquina (5 cifras significativas):
\[
    \pi = 3.14159265\ldots,
    \quad
    \fl(\pi) = 3.1416,
    \quad
    E_a = |3.14159265 - 3.1416| = 4.07 \times 10^{-5},
    \quad
    E_r \approx 1.30 \times 10^{-5}.
\]
\end{ejemplo}

\subsubsection{Importancia del error por redondeo}

\begin{observacion}
Los errores de redondeo pueden resultar especialmente críticos:
\begin{itemize}[leftmargin=1.5em]
    \item cuando intervienen cantidades muy grandes o muy pequeñas
          simultáneamente;
    \item cuando se restan números muy próximos (\emph{cancelación
          catastrófica});
    \item cuando un algoritmo repite muchas operaciones, acumulando pequeñas
          pérdidas de precisión (\emph{propagación de errores}).
\end{itemize}
\end{observacion}

\subsubsection{Cancelación catastrófica}

\begin{definicion}[Cancelación catastrófica]
La \textbf{cancelación catastrófica} se produce al restar dos números muy
próximos entre sí. Aunque ambos números estén bien representados en máquina,
el resultado puede perder muchas cifras significativas, amplificando el error
relativo de forma drástica.
\end{definicion}

\begin{ejemplo}
Considérese la resta exacta:
\[
    1\,436\,326\,440 - 1\,436\,326\,431 = 9.
\]
Ahora se incrementa el primer operando en un $0.001\%$:
\[
    1\,436\,326\,440 \times 1.00001 = 1\,436\,326\,440 + 14\,363
    \approx 1\,436\,340\,803.
\]
La nueva resta es:
\[
    1\,436\,340\,803 - 1\,436\,326\,431 = 14\,372.
\]

Una perturbación \emph{relativa} de $10^{-5}$ en el primer operando ha
producido una variación \emph{absoluta} de $14\,363$ en el resultado, que
pasó de 9 a $14\,372$. El error relativo en el resultado es enorme.

\medskip
\textbf{Análisis:} El número de cifras significativas perdidas es
$\lfloor \log_{10}(14372/9) \rfloor = \lfloor 3.2 \rfloor = 3$.
\end{ejemplo}

\begin{ejemplo}[Cancelación en fórmulas algebraicas]
Las raíces de $x^2 - bx + c = 0$ son
$x_{1,2} = \frac{b \pm \sqrt{b^2 - 4c}}{2}$.
Cuando $b \gg \sqrt{b^2 - 4c}$ (raíces muy distintas en magnitud), la raíz
pequeña sufre cancelación catastrófica. La solución estable consiste en
calcular la raíz grande con la fórmula estándar y obtener la pequeña mediante
la identidad $x_1 x_2 = c$:
\[
    x_1 = \frac{b + \sqrt{b^2 - 4c}}{2},
    \qquad
    x_2 = \frac{c}{x_1}.
\]
\end{ejemplo}

\subsubsection{Propagación de errores en operaciones básicas}

\begin{proposicion}[Propagación lineal del error]
Si $\tilde{x} = x + \delta x$ y $\tilde{y} = y + \delta y$ son aproximaciones
con errores $\delta x$ y $\delta y$, entonces:
\begin{align*}
    |\delta(x \pm y)| &\leq |\delta x| + |\delta y|
        \quad \text{(suma y resta)}, \\
    \frac{|\delta(xy)|}{|xy|} &\approx \frac{|\delta x|}{|x|}
        + \frac{|\delta y|}{|y|}
        \quad \text{(producto)}, \\
    \frac{|\delta(x/y)|}{|x/y|} &\approx \frac{|\delta x|}{|x|}
        + \frac{|\delta y|}{|y|}
        \quad \text{(cociente)}.
\end{align*}
En la resta, los errores absolutos se suman pero el resultado puede ser
pequeño, disparando el error relativo.
\end{proposicion}

\begin{demostracion}
Para el producto: $\tilde{x}\tilde{y} = (x+\delta x)(y+\delta y) = xy
+ x\,\delta y + y\,\delta x + \delta x\,\delta y$. Despreciando el término
cuadrático $\delta x\,\delta y$ (pequeño de segundo orden):
\[
    \delta(xy) \approx x\,\delta y + y\,\delta x
    \;\Rightarrow\;
    \frac{|\delta(xy)|}{|xy|} \lesssim \frac{|\delta x|}{|x|}
    + \frac{|\delta y|}{|y|}.
\]
El cociente se obtiene de forma análoga diferenciando $x/y$.
\end{demostracion}

\subsection{Relación entre error de truncamiento y error de redondeo}

\begin{observacion}
En la práctica, ambos tipos de error coexisten y sus efectos se contraponen:

\begin{itemize}[leftmargin=1.5em]
    \item Al \emph{aumentar} el orden de truncamiento (más términos de Taylor),
          el error de truncamiento disminuye, pero se realizan más operaciones
          aritméticas, acumulando más error de redondeo.
    \item Al \emph{disminuir} el paso $h$ en métodos de diferencias finitas,
          el error de truncamiento mejora como $O(h^p)$, pero el redondeo
          puede empeorar como $O(\epsilon_{\mathrm{máq}}/h)$.
\end{itemize}

Existe un \textbf{paso óptimo} $h^*$ que minimiza el error total. Para un
método de orden $p$ con constante $C$:
\[
    h^* \approx \left(\frac{p\,\epsilon_{\mathrm{máq}}}{C}\right)^{1/(p+1)}.
\]
\end{observacion}

\begin{cajaresultado}
\textbf{Conclusión del Tema 1.}
Todo método numérico debe analizarse desde dos perspectivas:
\[
    \text{precisión (error de truncamiento)}
    \quad \text{y} \quad
    \text{estabilidad (error de redondeo)}.
\]
El error de truncamiento mide la calidad de la aproximación teórica; el error
de redondeo mide la limitación de la aritmética finita. Ninguno de los dos
puede ignorarse en un análisis numérico riguroso.
\end{cajaresultado}

\subsection{Formulario del Tema 1}

\begin{center}
\begin{tabular}{ll}
\toprule
\textbf{Concepto} & \textbf{Fórmula} \\
\midrule
Error absoluto & $E_a = |x - \tilde{x}|$ \\[4pt]
Error relativo & $E_r = |x - \tilde{x}| / |x|$ \\[4pt]
Decimales correctos & $k = \lfloor -\log_{10}(2E_r) \rfloor$ \\[4pt]
Serie de Taylor & $f(x) = \sum_{k=0}^{n}
                  \frac{f^{(k)}(a)}{k!}(x-a)^k + R_n(x)$ \\[4pt]
Resto de Lagrange & $R_n(x) = \frac{f^{(n+1)}(\xi)}{(n+1)!}(x-a)^{n+1}$ \\[4pt]
Serie de Maclaurin & $f(x) = \sum_{k=0}^{n}
                     \frac{f^{(k)}(0)}{k!}x^k + R_n(x)$ \\[4pt]
$e^x$ (Maclaurin) & $e^x = \sum_{k=0}^{\infty} x^k / k!$ \\[4pt]
$\frac{1}{1+x^2}$ (Maclaurin) & $\sum_{k=0}^{\infty}(-1)^k x^{2k},\;|x|<1$ \\[4pt]
Error de redondeo & $E_{\mathrm{red}} = |x - \fl(x)|$ \\[4pt]
Épsilon de máquina (IEEE 754) & $\epsilon_{\mathrm{máq}} \approx 2.22\times
                                 10^{-16}$ \\
\bottomrule
\end{tabular}
\end{center}

\newpage

% ==================================================
% TEMA 2
% ==================================================
\section{Tema 2: Ecuaciones no lineales}

\subsection{Introducción}

Muchas ecuaciones de interés científico y técnico no pueden resolverse de forma
explícita. En tales casos se busca una aproximación numérica de una raíz de la
ecuación
\[
f(x)=0.
\]

A lo largo de este tema se desarrollan distintos enfoques: detección gráfica de
raíces, métodos cerrados basados en intervalos, métodos abiertos de iteración,
y procedimientos especiales para polinomios con coeficientes enteros.

\subsection{Planteamiento general}

\begin{definicion}[Raíz de una ecuación no lineal]
Un número $r\in\R$ es una raíz de $f(x)=0$ si
\[
f(r)=0.
\]
\end{definicion}

\begin{definicion}[Problema numérico]
Dada una función $f:\R\to\R$, se desea construir una sucesión
\[
\{x_n\}_{n\ge 0}
\]
tal que
\[
x_n \to r, \qquad f(r)=0.
\]
\end{definicion}

\subsection{Tipos de ecuaciones no lineales}

Las ecuaciones no lineales pueden involucrar polinomios, funciones
trascendentes, funciones racionales, funciones exponenciales, trigonométricas
o combinaciones de ellas. Ejemplos típicos son:
\[
x^3-x-1=0,\qquad \cos x - x =0,\qquad e^{-x}-x=0,\qquad
\tan x - x -0.5=0.
\]

\subsection{Métodos iterativos}

\begin{definicion}[Método iterativo]
Un método iterativo genera una sucesión
\[
x_{n+1}=\Phi(x_n)
\]
a partir de un valor inicial $x_0$. La esperanza es que esta sucesión converja
a una raíz.
\end{definicion}

\begin{observacion}
Los métodos iterativos se clasifican en:
\begin{itemize}[leftmargin=1.5em]
    \item \textbf{métodos cerrados}, si mantienen un intervalo que contiene la
          raíz;
    \item \textbf{métodos abiertos}, si utilizan uno o varios puntos iniciales
          sin garantía automática de encierro.
\end{itemize}
\end{observacion}

\subsection{Métodos gráficos}

\subsubsection{Idea geométrica}

El método gráfico consiste en representar la función $y=f(x)$ y localizar
visualmente los puntos donde corta al eje $x$.

\begin{definicion}[Aislamiento gráfico de raíces]
Se dice que una raíz está aislada en un intervalo $[a,b]$ si la representación
de $f$ sugiere que existe un único corte con el eje $x$ dentro de ese
intervalo.
\end{definicion}

\begin{observacion}
El método gráfico no proporciona una raíz precisa, pero sí permite:
\begin{itemize}[leftmargin=1.5em]
    \item detectar la existencia de raíces;
    \item estimar intervalos iniciales;
    \item elegir valores adecuados para métodos posteriores.
\end{itemize}
\end{observacion}

\begin{figure}[H]
\centering
\begin{tikzpicture}[scale=0.95]
    \draw[->] (-3.2,0) -- (3.4,0) node[right] {$x$};
    \draw[->] (0,-2.2) -- (0,2.8) node[above] {$y$};
    \draw[blue!70!black, thick, domain=-2.8:2.8, samples=200]
        plot (\x,{0.25*(\x)^3 - \x + 0.2});
    \draw[dashed] (0.2,-2) -- (0.2,2.4);
    \filldraw (0.2,0) circle (2pt) node[below right] {$r$};
    \node at (2.35,2.1) {$y=f(x)$};
\end{tikzpicture}
\caption{Visualización gráfica de una raíz aproximada.}
\end{figure}

\subsection{Método de bisección}

\subsubsection{Fundamento teórico}

\begin{teorema}[Bolzano]
Sea $f$ continua en $[a,b]$. Si
\[
f(a)f(b)<0,
\]
entonces existe al menos un $r\in(a,b)$ tal que
\[
f(r)=0.
\]
\end{teorema}

\begin{demostracion}
Es una consecuencia directa del teorema del valor intermedio. Como $f(a)$ y
$f(b)$ tienen signos opuestos y $f$ es continua, la función debe anularse en
algún punto intermedio.
\end{demostracion}

\subsubsection{Algoritmo}

Dado un intervalo inicial $[a_0,b_0]$ con $f(a_0)f(b_0)<0$, se define el
punto medio
\[
x_n=\frac{a_n+b_n}{2}.
\]
Luego:
\[
\text{si } f(a_n)f(x_n)<0,\quad [a_{n+1},b_{n+1}]=[a_n,x_n],
\]
\[
\text{si } f(x_n)f(b_n)<0,\quad [a_{n+1},b_{n+1}]=[x_n,b_n].
\]

\begin{cajaresultado}
\textbf{Fórmula de bisección:}
\[
x_n=\frac{a_n+b_n}{2}.
\]
\end{cajaresultado}

\subsubsection{Cota del error}

\begin{proposicion}
Después de $n$ iteraciones, la longitud del intervalo es
\[
b_n-a_n=\frac{b_0-a_0}{2^n}.
\]
Por tanto,
\[
|r-x_n|\le \frac{b_0-a_0}{2^{n+1}}.
\]
\end{proposicion}

\begin{demostracion}
Cada iteración reduce el intervalo a la mitad. Tras $n$ pasos:
\[
b_n-a_n=\frac{b_0-a_0}{2^n}.
\]
Como $x_n$ es el punto medio del intervalo que contiene la raíz, la distancia
máxima entre $x_n$ y $r$ es la mitad de la longitud del intervalo.
\end{demostracion}

\subsubsection{Criterios de parada}

Se suele detener el método cuando se cumple alguna de estas condiciones:
\[
|f(x_n)|<\eps, \qquad |x_n-x_{n-1}|<\eps, \qquad \frac{b_n-a_n}{2}<\eps.
\]

\subsubsection{Ejemplo}

\begin{ejemplo}
Aproximar una raíz de $f(x)=x^3-x-2$ en el intervalo $[1,2]$.

Como $f(1)=-2$ y $f(2)=4$, hay una raíz en $(1,2)$.

Primera iteración:
$x_1=1.5$, $f(1.5)=-0.125$. La raíz sigue en $[1.5,2]$.

Segunda iteración:
$x_2=1.75$, $f(1.75)=1.609375$. Raíz en $[1.5,1.75]$.

Tercera iteración:
$x_3=1.625$, $f(1.625)=0.666015625$. Raíz en $[1.5,1.625]$.

Repitiendo el proceso se obtiene una aproximación convergente a la raíz real.
\end{ejemplo}

\subsection{Método de regla falsa}

\subsubsection{Idea geométrica}

La regla falsa o \emph{regula falsi} mejora la bisección sustituyendo el punto
medio por la intersección con el eje $x$ de la secante que une los extremos del
intervalo.

\begin{proposicion}
El corte con el eje $x$ viene dado por
\[
x_n = \frac{a_n f(b_n)-b_n f(a_n)}{f(b_n)-f(a_n)}.
\]
\end{proposicion}

\begin{demostracion}
La ecuación de la recta secante puede escribirse como
\[
y-f(a_n)=\frac{f(b_n)-f(a_n)}{b_n-a_n}(x-a_n).
\]
Imponiendo $y=0$ y despejando $x$ se obtiene la expresión indicada.
\end{demostracion}

\begin{cajaresultado}
\textbf{Fórmula de la regla falsa:}
\[
x_n = \frac{a_n f(b_n)-b_n f(a_n)}{f(b_n)-f(a_n)}.
\]
\end{cajaresultado}

\subsubsection{Actualización del intervalo}

Se conserva el extremo que mantenga el cambio de signo:
\[
\text{si } f(a_n)f(x_n)<0,\quad [a_{n+1},b_{n+1}]=[a_n,x_n],
\]
\[
\text{si } f(x_n)f(b_n)<0,\quad [a_{n+1},b_{n+1}]=[x_n,b_n].
\]

\subsubsection{Observación}

La regla falsa suele converger más deprisa que la bisección cuando la función
es aproximadamente lineal en el intervalo. Sin embargo, en algunos problemas uno
de los extremos permanece fijo durante muchas iteraciones, ralentizando el
proceso.

\subsection{Método de punto fijo}

\subsubsection{Planteamiento}

Considérese la ecuación $f(x)=0$. Si puede reescribirse en la forma $x=g(x)$,
entonces una raíz de la ecuación es un punto fijo de $g$.

\begin{definicion}[Punto fijo]
Un número $p$ es punto fijo de $g$ si $g(p)=p$.
\end{definicion}

El método genera la sucesión $x_{n+1}=g(x_n)$.

\subsubsection{Interpretación geométrica}

El método busca el punto de intersección entre las curvas $y=x$ e $y=g(x)$.

\begin{figure}[H]
\centering
\begin{tikzpicture}[scale=1.0]
    \draw[->] (-0.5,0) -- (4.5,0) node[right] {$x$};
    \draw[->] (0,-0.5) -- (0,4.3) node[above] {$y$};
    \draw[thick] (0,0) -- (4,4) node[right] {$y=x$};
    \draw[blue!70!black, thick, domain=0.2:4, samples=100]
        plot (\x,{3.2/(\x+0.5)+0.35});
    \filldraw (1.4,1.4) circle (2pt) node[below right] {$p$};
\end{tikzpicture}
\caption{Interpretación geométrica del método de punto fijo.}
\end{figure}

\subsubsection{Teorema de convergencia}

\begin{teorema}[Convergencia local del método de punto fijo]
Sea $g\in C^1([a,b])$ tal que $g([a,b])\subseteq [a,b]$ y exista una constante
$k<1$ con $|g'(x)|\le k$ para todo $x\in[a,b]$. Entonces:
\begin{enumerate}[label=\roman*)]
    \item existe un único punto fijo $p\in[a,b]$;
    \item para cualquier $x_0\in[a,b]$, la sucesión $x_{n+1}=g(x_n)$
          converge a $p$.
\end{enumerate}
\end{teorema}

\begin{observacion}
La condición $|g'(x)|<1$ garantiza una contracción. Si en cambio $|g'(p)|>1$,
el punto fijo tiende a ser repulsivo y el método suele divergir.
\end{observacion}

\subsubsection{Ejemplo clásico}

\begin{ejemplo}
Resolver aproximadamente $\cos x - x = 0$.
La ecuación puede escribirse como $x=\cos x$, con $g(x)=\cos x$.
Como $|g'(x)|=|\!-\!\sin x|\le \sin 1 < 1$ en $[0,1]$, el método converge.

Con $x_0=0.5$:
\[
x_1\approx 0.87758256,\quad x_2\approx 0.63901249,\quad
x_3\approx 0.80268510,\quad x_4\approx 0.69477803.
\]
La sucesión converge a $p\approx 0.73908513$.
\end{ejemplo}

\subsection{Método de Newton-Raphson}

\subsubsection{Idea geométrica}

El método de Newton aproxima la función por su recta tangente en un punto $x_n$
y toma como nuevo valor el corte de esa tangente con el eje $x$.

\begin{figure}[H]
\centering
\begin{tikzpicture}[scale=0.95]
    \draw[->] (-0.5,0) -- (4.7,0) node[right] {$x$};
    \draw[->] (0,-0.5) -- (0,4.2) node[above] {$y$};
    \draw[blue!70!black, thick, domain=0.4:4.2, samples=150]
        plot (\x,{0.18*(\x-1.2)^3 + 0.85*(\x-1.2)+0.45});
    \draw[dashed] (3.0,0) -- (3.0,2.12);
    \filldraw (3.0,2.12) circle (2pt) node[above right] {$(x_n,f(x_n))$};
    \draw[red!70!black, thick] (1.7,-0.1) -- (4.0,3.5);
    \filldraw (1.85,0) circle (2pt) node[below] {$x_{n+1}$};
\end{tikzpicture}
\caption{Interpretación geométrica del método de Newton-Raphson.}
\end{figure}

\subsubsection{Deducción}

La recta tangente a $y=f(x)$ en $x_n$ es
$y-f(x_n)=f'(x_n)(x-x_n)$.
Imponiendo $y=0$:

\begin{cajaresultado}
\textbf{Fórmula de Newton-Raphson:}
\[
x_{n+1}=x_n-\frac{f(x_n)}{f'(x_n)}.
\]
\end{cajaresultado}

\subsubsection{Convergencia local}

\begin{teorema}[Convergencia cuadrática]
Sea $r$ una raíz simple de $f$, con $f\in C^2$ en un entorno de $r$, y
supóngase $f(r)=0$, $f'(r)\neq 0$. Si $x_0$ está suficientemente cerca de $r$,
entonces la sucesión de Newton converge a $r$ y el error satisface
\[
|e_{n+1}| \approx C |e_n|^2, \qquad e_n=x_n-r.
\]
\end{teorema}

\begin{observacion}
La convergencia cuadrática significa que, una vez cerca de la raíz, el número
de cifras correctas aproximadamente se duplica en cada iteración.
\end{observacion}

\subsubsection{Ejemplo: raíz de $\cos x-x=0$}

\begin{ejemplo}
Para $f(x)=\cos x - x$, $f'(x)=-\sin x -1$:
\[
x_{n+1}=x_n+\frac{\cos x_n - x_n}{\sin x_n +1}.
\]
Con $x_0=1$: $x_1\approx 0.75036387$, $x_2\approx 0.73911289$,
$x_3\approx 0.73908513$.
\end{ejemplo}

\subsubsection{Ejemplo: aproximación de $\sqrt{2}$}

\begin{ejemplo}
Para $f(x)=x^2-2$, $f'(x)=2x$:
\[
x_{n+1}=\frac{x_n^2+2}{2x_n}.
\]
Con $x_0=1$: $x_1=1.5$, $x_2=1.41\overline{6}$, $x_3=1.41421569$,
$x_4=1.41421356$.
\end{ejemplo}

\subsection{Método de Newton-Raphson mejorado}

\subsubsection{Motivación}

Cuando la raíz tiene multiplicidad mayor que uno, el método de Newton clásico
pierde su convergencia cuadrática y pasa a converger linealmente.

\begin{definicion}[Raíz múltiple]
Se dice que $r$ es raíz de multiplicidad $m$ de $f$ si
\[
f(r)=f'(r)=\cdots=f^{(m-1)}(r)=0, \qquad f^{(m)}(r)\neq 0.
\]
\end{definicion}

\begin{cajaresultado}
\textbf{Newton mejorado para raíces múltiples:}
\[
x_{n+1}=x_n-\frac{f(x_n)f'(x_n)}{[f'(x_n)]^2-f(x_n)f''(x_n)}.
\]
\end{cajaresultado}

\subsection{Método de la secante}

\begin{cajaresultado}
\textbf{Fórmula de la secante:}
\[
x_{n+1} = x_n - \frac{f(x_n)(x_n-x_{n-1})}{f(x_n)-f(x_{n-1})}.
\]
\end{cajaresultado}

\begin{observacion}
El método de la secante requiere dos valores iniciales $x_0$ y $x_1$. Su
convergencia suele ser más rápida que la de punto fijo y comparable a Newton,
aunque algo más lenta que la cuadrática.
\end{observacion}

\subsection{Raíces de polinomios con coeficientes enteros}

\subsubsection{Teorema de la raíz racional}

\begin{teorema}[Raíz racional]
Sea $P(x)=a_nx^n+a_{n-1}x^{n-1}+\cdots+a_0$ con coeficientes enteros. Si
$p/q$ (irreducible) es raíz racional de $P$, entonces $p\mid a_0$ y
$q\mid a_n$.
\end{teorema}

\begin{ejemplo}
Para $P(x)=2x^3+4x^2-22x-24$, las raíces racionales candidatas son
$\pm 1,\pm 2,\pm 3,\pm 4,\pm 6,\pm 8,\pm 12,\pm 24,\pm\tfrac12,\pm\tfrac32$.
Probando: $P(-4)=P(-1)=P(3)=0$.
\end{ejemplo}

\subsubsection{Regla de los signos de Descartes}

\begin{teorema}[Regla de Descartes]
El número de raíces positivas de un polinomio real es igual al número de
cambios de signo de $P(x)$ o menor que él en un número par. El número de
raíces negativas se obtiene aplicando la misma regla a $P(-x)$.
\end{teorema}

\subsection{División sintética}

La división sintética permite dividir rápidamente $P(x)$ entre un factor
lineal $x-c$. Si $P(c)=0$, entonces $x-c$ es factor y la división permite
rebajar el grado del polinomio.

\begin{ejemplo}
Sea $P(x)=x^3+5x^2+x-10$. Como $P(-2)=0$:
\[
P(x)=(x+2)(x^2+3x-5),
\qquad
x=\frac{-3\pm\sqrt{29}}{2}.
\]
\end{ejemplo}

\begin{cajaresultado}
\textbf{Resumen del Tema 2.}
\begin{center}
\begin{tabular}{p{3.2cm}p{3.2cm}p{5.3cm}}
\toprule
Método & Tipo & Idea básica \\
\midrule
Gráfico & preliminar & detectar intervalos con raíces \\
Bisección & cerrado & dividir el intervalo por mitades \\
Regla falsa & cerrado & usar la secante entre extremos \\
Punto fijo & abierto & iterar $x_{n+1}=g(x_n)$ \\
Newton-Raphson & abierto & usar la tangente \\
Newton mejorado & abierto & corregir multiplicidad \\
Secante & abierto & aproximar la derivada \\
Raíz racional & algebraico & listar candidatas racionales \\
División sintética & algebraico & factorizar y reducir grado \\
\bottomrule
\end{tabular}
\end{center}
\end{cajaresultado}

\newpage

% ==================================================
% TEMA 3
% ==================================================
\section{Tema 3: Interpolación y aproximación de funciones}

\subsection{Introducción}

En muchos problemas no se conoce la función exacta, sino sólo un conjunto de
datos experimentales o tabulados. En tal situación surgen dos tareas
diferentes:

\begin{itemize}[leftmargin=1.5em]
    \item \textbf{Interpolar}: construir una función que pase exactamente por
          los datos.
    \item \textbf{Aproximar}: construir una función que describa la tendencia
          global de los datos, aunque no pase exactamente por todos ellos.
\end{itemize}

Este tema desarrolla ambos enfoques: ajuste por mínimos cuadrados y métodos
de interpolación polinómica.

\subsection{Método de mínimos cuadrados}

\subsubsection{Planteamiento general}

Sea un conjunto de datos $\{(x_1,y_1),\dots,(x_n,y_n)\}$.
Se desea ajustar un modelo $y=f(x;\beta_0,\beta_1,\dots,\beta_m)$.

\begin{definicion}[Error cuadrático total]
\[
E(\beta_0,\dots,\beta_m)=\sum_{i=1}^n \left(y_i-f(x_i;\beta_0,\dots,\beta_m)
\right)^2.
\]
\end{definicion}

\begin{cajaresultado}
\textbf{Principio de mínimos cuadrados:} los mejores parámetros minimizan
$E = \sum_{i=1}^n e_i^2$.
\end{cajaresultado}

\subsection{Regresión lineal}

El modelo lineal es $y=\beta_0+\beta_1 x$.

\begin{proposicion}
Las ecuaciones normales de la regresión lineal son
\[
n\beta_0+\beta_1\sum x_i = \sum y_i,
\qquad
\beta_0\sum x_i + \beta_1\sum x_i^2 = \sum x_i y_i.
\]
\end{proposicion}

\begin{proposicion}
Si $\Delta = n\sum x_i^2-\left(\sum x_i\right)^2 \neq 0$, entonces
\[
\beta_1=\frac{n\sum x_i y_i-\left(\sum x_i\right)\left(\sum y_i\right)}
             {\Delta},
\qquad
\beta_0=\frac{\left(\sum x_i^2\right)\left(\sum y_i\right)
             -\left(\sum x_i y_i\right)\left(\sum x_i\right)}{\Delta}.
\]
\end{proposicion}

\subsubsection{Ejemplo detallado}

\begin{ejemplo}
Ajustar por regresión lineal los puntos
$(1,0.5),(2,2.5),(3,2),(4,4),(5,3.5),(6,6),(7,5.5)$.

\begin{center}
\begin{tabular}{ccccc}
\toprule
$i$ & $x_i$ & $y_i$ & $x_i y_i$ & $x_i^2$ \\
\midrule
1 & 1 & 0.5 & 0.5 & 1 \\
2 & 2 & 2.5 & 5.0 & 4 \\
3 & 3 & 2.0 & 6.0 & 9 \\
4 & 4 & 4.0 & 16.0 & 16 \\
5 & 5 & 3.5 & 17.5 & 25 \\
6 & 6 & 6.0 & 36.0 & 36 \\
7 & 7 & 5.5 & 38.5 & 49 \\
\midrule
$\sum$ & 28 & 24 & 119.5 & 140 \\
\bottomrule
\end{tabular}
\end{center}

$\beta_1 = \frac{7(119.5)-28(24)}{7(140)-28^2} = \frac{164.5}{196}
\approx 0.8393$,\quad
$\beta_0 = \frac{140(24)-119.5(28)}{196} = \frac{14}{196} \approx 0.0714$.

Recta ajustada: $\hat{y} = 0.0714 + 0.8393\,x$.
\end{ejemplo}

\subsection{Regresión polinomial}

Para un polinomio de grado 2: $y=\beta_0+\beta_1 x+\beta_2 x^2$.

\begin{proposicion}
Las ecuaciones normales del ajuste cuadrático son
\[
n\beta_0+\left(\sum x_i\right)\beta_1+\left(\sum x_i^2\right)\beta_2
= \sum y_i,
\]
\[
\left(\sum x_i\right)\beta_0+\left(\sum x_i^2\right)\beta_1
+\left(\sum x_i^3\right)\beta_2 = \sum x_i y_i,
\]
\[
\left(\sum x_i^2\right)\beta_0+\left(\sum x_i^3\right)\beta_1
+\left(\sum x_i^4\right)\beta_2 = \sum x_i^2 y_i.
\]
\end{proposicion}

\subsection{Mínimos cuadrados para un modelo exponencial}

\begin{proposicion}
El modelo $y=Ae^{Bx}$ se linealiza como $Y = \alpha + Bx$ con
$Y=\ln y$, $\alpha=\ln A$. Una vez estimados $\alpha$ y $B$: $A=e^\alpha$.
\end{proposicion}

\subsection{Linealización}

\begin{center}
\begin{tabular}{ccc}
\toprule
Modelo original & Transformación & Modelo lineal \\
\midrule
$y=Ae^{Bx}$ & $Y=\ln y$ & $Y=\ln A + Bx$ \\
$y=Ax^B$ & $Y=\ln y,\ X=\ln x$ & $Y=\ln A + BX$ \\
$y=\frac{1}{A+Bx}$ & $Y=\frac{1}{y}$ & $Y=A+Bx$ \\
\bottomrule
\end{tabular}
\end{center}

\subsection{Interpolación de Newton}

\begin{definicion}[Diferencia dividida]
\[
f[x_i,x_{i+1}] = \frac{f(x_{i+1})-f(x_i)}{x_{i+1}-x_i},
\qquad
f[x_i,\dots,x_{i+k}]
= \frac{f[x_{i+1},\dots,x_{i+k}]-f[x_i,\dots,x_{i+k-1}]}{x_{i+k}-x_i}.
\]
\end{definicion}

\begin{teorema}[Polinomio de Newton]
\[
P_n(x)=f[x_0]+f[x_0,x_1](x-x_0)+f[x_0,x_1,x_2](x-x_0)(x-x_1)+\cdots
\]
\[
\hspace{2.2cm}
+f[x_0,\dots,x_n](x-x_0)(x-x_1)\cdots(x-x_{n-1}).
\]
\end{teorema}

\begin{ejemplo}
Interpolar $(1,1),(2,4),(3,9)$.
Diferencias: $f[1,2]=3$, $f[2,3]=5$, $f[1,2,3]=1$.
\[
P_2(x)=1+3(x-1)+1(x-1)(x-2)=x^2.
\]
\end{ejemplo}

\subsection{Interpolación de Lagrange}

\begin{definicion}[Polinomio básico de Lagrange]
\[
L_i(x)=\prod_{\substack{j=0\\ j\neq i}}^n \frac{x-x_j}{x_i-x_j},
\qquad L_i(x_j)=\delta_{ij}.
\]
\end{definicion}

\begin{teorema}[Polinomio de Lagrange]
\[
P_n(x)=\sum_{i=0}^n y_i\,L_i(x).
\]
\end{teorema}

\begin{ejemplo}
Interpolar $(0,1),(1,3),(2,2)$:
\[
P_2(x) = -\frac{3}{2}x^2+\frac{7}{2}x+1.
\]
\end{ejemplo}

\subsection{Error de interpolación}

\begin{teorema}[Resto de interpolación]
Sea $f\in C^{n+1}[a,b]$ y $P_n$ su polinomio interpolador. Entonces existe
$\xi$ entre los nodos y $x$ tal que
\[
f(x)-P_n(x)=\frac{f^{(n+1)}(\xi)}{(n+1)!}\prod_{i=0}^n (x-x_i).
\]
\end{teorema}

\subsection{Comparación entre interpolación y aproximación}

\begin{center}
\begin{tabular}{p{4cm}p{5cm}p{5cm}}
\toprule
Aspecto & Interpolación & Aproximación por mínimos cuadrados \\
\midrule
Relación con los datos & pasa exactamente por todos & no necesariamente \\
Objetivo & reproducir los valores dados & captar la tendencia global \\
Uso típico & tablas, nodos exactos & datos experimentales con ruido \\
Error & nulo en los nodos & minimiza suma de cuadrados \\
\bottomrule
\end{tabular}
\end{center}

\begin{cajaresultado}
$\text{Interpolar} \neq \text{Ajustar}$.
La interpolación reproduce exactamente los datos; el ajuste por mínimos
cuadrados busca el modelo que mejor explica el conjunto de observaciones en
sentido cuadrático promedio.
\end{cajaresultado}

\newpage

% ==================================================
% FORMULARIO GENERAL
% ==================================================
\section{Formulario general}

\subsection{Errores (Tema 1)}

\[
E_a = |x-\tilde{x}|,
\qquad
E_r = \frac{|x-\tilde{x}|}{|x|},
\qquad
k \approx \left\lfloor -\log_{10}(2E_r)\right\rfloor.
\]

\subsection{Taylor y Maclaurin}

\[
f(x)=\sum_{k=0}^{n}\frac{f^{(k)}(a)}{k!}(x-a)^k+R_n(x),
\qquad
R_n(x)=\frac{f^{(n+1)}(\xi)}{(n+1)!}(x-a)^{n+1}.
\]

\[
e^x=\sum_{k=0}^\infty \frac{x^k}{k!},
\qquad
\frac{1}{1+x^2}=\sum_{k=0}^\infty (-1)^k x^{2k},\quad |x|<1.
\]

\subsection{Ecuaciones no lineales}

\[
\text{Bisección:}\quad x_n=\frac{a_n+b_n}{2},
\qquad
|r-x_n|\le \frac{b_0-a_0}{2^{n+1}}.
\]

\[
\text{Regla falsa:}\quad
x_n=\frac{a_n f(b_n)-b_n f(a_n)}{f(b_n)-f(a_n)}.
\]

\[
\text{Punto fijo:}\quad x_{n+1}=g(x_n).
\qquad
\text{Newton-Raphson:}\quad
x_{n+1}=x_n-\frac{f(x_n)}{f'(x_n)}.
\]

\[
\text{Newton mejorado:}\quad
x_{n+1}=x_n-\frac{f(x_n)f'(x_n)}{[f'(x_n)]^2-f(x_n)f''(x_n)}.
\]

\[
\text{Secante:}\quad
x_{n+1}=x_n-\frac{f(x_n)(x_n-x_{n-1})}{f(x_n)-f(x_{n-1})}.
\]

\subsection{Mínimos cuadrados}

\[
E=\sum_{i=1}^n \left(y_i-f(x_i;\beta)\right)^2,
\qquad
\beta_1=\frac{n\sum x_i y_i-\left(\sum x_i\right)\left(\sum y_i\right)}
             {n\sum x_i^2-\left(\sum x_i\right)^2}.
\]

\[
\beta_0=\frac{\left(\sum x_i^2\right)\left(\sum y_i\right)
             -\left(\sum x_i y_i\right)\left(\sum x_i\right)}
             {n\sum x_i^2-\left(\sum x_i\right)^2}.
\]

\subsection{Interpolación}

\[
P_n(x)=f[x_0]+f[x_0,x_1](x-x_0)+\cdots
       +f[x_0,\dots,x_n]\prod_{j=0}^{n-1}(x-x_j).
\]

\[
L_i(x)=\prod_{\substack{j=0\\j\neq i}}^n \frac{x-x_j}{x_i-x_j},
\qquad
P_n(x)=\sum_{i=0}^n y_i L_i(x).
\]

\[
f(x)-P_n(x)=\frac{f^{(n+1)}(\xi)}{(n+1)!}\prod_{i=0}^n (x-x_i).
\]

\end{document}